\section{Sampling the integral cut comparison}
\label{sec:sampling}

The comparison in \eqref{eq:44} allows arbitrary bounded functions at
each single site.  To use it against the labels in \eqref{eq:19}, we
must also allow the testing signs to depend on the entire sampled
matrix.  We do this by approximating an optimizing row-sign vector
using a small set of columns, and then applying concentration
simultaneously to the resulting finite family of optimizations.

The small-column approximation is related to the sampling argument of
Alon, Fern\'andez de la Vega, Kannan, and Karpinski
\citep[Lemma~3]{AlonEtAl2003Sampling} and to the proof of
\citet[Theorem~4.6]{BorgsEtAl2008}. Their arguments approximate optimizing
cuts by a small sample and then control a finite family of tests.
Here the kernels depend on $B$ and on position, and are controlled by
degree bounds and integrated row-second-moment bounds. We therefore
prove the required sampling transfer directly rather than invoking those
results as a black box.

Throughout this section the regularity parameters, $C_*$, $\eta$,
$C_6$, $\epsilon_1$, and the coarse partition are fixed.  In particular,
$T=\lfloor B^{0.32}\rfloor$ and $M=\lceil C_6T\rceil\asymp B^{0.32}$.
All limits and $o(1)$ terms in the probabilistic argument refer to
$B\to\infty$ with these parameters fixed.  We work at one value of $s$
in the fixed compact range of the smooth weights.  Every bound below
is uniform in that value of $s$; no simultaneous event over all $s$
will be needed.

Let $\mathbf S=(S_1,\ldots,S_M)$ have the independent site law from
Section~\ref{sec:graph}.  Write
\[
 \mathcal E_{ik}(S_i,S_k;s)
   =\mathcal L_{ik}(S_i,S_k;s)-\mathcal M_{ik}(S_i,S_k;s),
 \qquad
 E_0=C e^{-c_3C_*}+C_\eta\epsilon_1.
\]
Here the fixed constants in $E_0$ are large enough for
\eqref{eq:44}.  As established there, the coefficient $C$ of
$e^{-c_3C_*}$ can be chosen independently of $\eta$ and $C_6$.
The matrices are real and symmetric, and are zero on the diagonal
and on all nonallowed pairs.  We retain the normalization
\[
 \|H\|_\square
 =\frac1M\max_{\epsilon,\delta\in\{-1,1\}^M}
       \left|\sum_{i,k}H_{ik}\epsilon_i\delta_k\right|.
\]

\begin{proposition}[Sampling the type-kernel comparison]
\label{prop:sampling-comparison}
For every fixed $\epsilon_2>0$,
\begin{equation}
 \E\big\|\big(\mathcal E_{ik}(S_i,S_k;s)\big)\big\|_\square
       \le E_0+\epsilon_2+o(1),
 \label{eq:45}
\end{equation}
uniformly for $s$ in the fixed compact range.
\end{proposition}

We first record the degree and integrability estimates required in
the proof.  Use the nonnegative symmetric envelope
\[
 \mathcal H_{ik}=\mathcal L_{ik}+|\mathcal M_{ik}|,
 \qquad |\mathcal E_{ik}|\le\mathcal H_{ik}.
\]
Dependence on the endpoint types and on $s$ is suppressed when no
confusion is possible.  By the every-type bound \eqref{eq:23} and
the deterministic absolute row bounds for \eqref{eq:43}, there is
a constant $C_0$ such that
\[
 \E\left[\sum_k\mathcal H_{ik}\,\middle|\,S_i\right]\le C_0
       \qquad\text{for every row type }S_i.
\]
By \eqref{eq:24} and the squared-row bounds following
\eqref{eq:44}, there is a constant $C_1$ such that
\[
 \sum_k\E\mathcal H_{ik}^2\le q_B,
 \qquad
 \sum_k\E\mathcal E_{ik}^2\le q_B,
 \qquad q_B=C_1B^{-0.21},
       \qquad 1\le i\le M.
\]
Indeed $(\mathcal L_{ik}+|\mathcal M_{ik}|)^2
\le2\mathcal L_{ik}^2+2\mathcal M_{ik}^2$, and
$1/T=O(B^{-0.32})$.  These squared-row bounds are integrated over
the row type; no uniform conditional second-moment assertion is
being made.

\begin{lemma}[A common degree event and uniform integrability]
\label{lem:sampling-good}
There is a fixed $K>0$ such that, on an event $\mathcal G_B$ of
probability $1-O(B^{-0.01})-O(B^{-0.03})$,
\begin{equation}
 \frac1M\sum_{i,k}\mathcal E_{ik}^2\le B^{-0.20},
 \qquad
 \max_i\sum_k|\mathcal E_{ik}|\le K B^{0.07}.
 \label{eq:46}
\end{equation}
The event may also be required to satisfy
$\max_i\sum_k\mathcal H_{ik}\le K B^{0.07}$.
Moreover, the normalized total absolute mass
\[
 A_B=\frac1M\sum_{i,k}|\mathcal E_{ik}|
\]
has bounded second moment, uniformly in $B$ and $s$.
\end{lemma}

\begin{proof}
Put $D_i=\sum_k\mathcal H_{ik}$ and $d_B=K B^{0.07}$,
where $K$ is any fixed sufficiently large constant.  Conditional
on $S_i$, the summands of $D_i$ are independent, since each
off-diagonal summand depends on one other independent site.
Consequently
\[
 \E\Var(D_i\mid S_i)
 \le\sum_k\E\mathcal H_{ik}^2\le q_B,
 \qquad
 \E D_i^2\le C_0^2+q_B.
\]
For sufficiently large $B$, $d_B\ge2C_0$.  Conditional Chebyshev
followed by averaging the row type gives
\[
 \Prob(D_i>d_B)
 \le\frac{q_B}{(d_B-C_0)^2}
 \le\frac{4q_B}{d_B^2}.
\]
Thus the event
\[
 \Omega_{\rm deg}
   =\{\max_iD_i\le d_B\}
\]
has complement of probability
$O(Mq_B/d_B^2)=O(B^{-0.03})$.
Markov's inequality gives
\[
 \Prob\left(\frac1M\sum_{i,k}\mathcal E_{ik}^2>B^{-0.20}\right)
       \le q_B B^{0.20}=O(B^{-0.01}).
\]
Take
\[
 \mathcal G_B=\Omega_{\rm deg}\cap
   \left\{\frac1M\sum_{i,k}\mathcal E_{ik}^2\le B^{-0.20}\right\}.
\]
This implies \eqref{eq:46}.
Finally, Jensen's inequality across rows gives
\[
 \E A_B^2
 \le\E\left(\frac1M\sum_iD_i\right)^2
 \le\frac1M\sum_i\E D_i^2
 \le C_0^2+q_B.
\]
In particular, for any events $\mathcal B_B$ with probabilities
tending to zero uniformly in $s$,
$\E[A_B\ind_{\mathcal B_B}]=o(1)$ by Cauchy--Schwarz.
\end{proof}

The same degree estimate applies to the graph induced by any fixed
subset of indices: deleting terms decreases both its conditional
expected degrees and its integrated squared-row sums.  In particular,
the probability of a degree exceeding $d_B$ is still
$O(B^{-0.03})$ after a union over its at most $M$ rows.  This remains
true conditional on the types at indices outside the subset, since
the remaining types retain their original independent laws.

\begin{lemma}[Approximation by a small set of columns]
\label{lem:sampling-skeleton}
Put $m_2=\lfloor M B^{-0.18}\rfloor$.  On $\mathcal G_B$ there
are a subset $J'\subset\{1,\ldots,M\}$ of size $m_2$ and
signs $(\delta_j)_{j\in J'}$ such that, with
\[
 \widetilde g_i
 =\sgn\left(\sum_{j\in J'}\mathcal E_{ij}\delta_j\right),
 \qquad R' =\{1,\ldots,M\}\setminus J',
\]
and assigning the sign $+1$ at zero for all sign choices in this section,
\[
 \|\mathcal E\|_\square
 \le\frac1M\sum_{k\in R'}
       \left|\sum_{i\in R'}\mathcal E_{ik}\widetilde g_i\right|
       +O(B^{-0.01})+O(B^{-0.11}).
\]
The constants are independent of the realized matrix on $\mathcal G_B$.
\end{lemma}

\begin{proof}
Fix the realized matrix.  Choose column signs $\delta_k$ attaining
the cut norm and choose the best-response row signs, so the
unnormalized maximum is the positive quantity
$\sum_i|r_i|$, where $r_i=\sum_k\mathcal E_{ik}\delta_k$.
For this deterministic matrix only, sample a uniform subset $J'$
of $m_2$ columns and form the unbiased estimate
\[
 \widehat r_i=\frac{M}{m_2}\sum_{j\in J'}\mathcal E_{ij}\delta_j.
\]
The variance formula for sampling without replacement yields
\[
 \E_{J'}|\widehat r_i-r_i|^2
       \le\frac{M}{m_2}\sum_k\mathcal E_{ik}^2.
\]
For real numbers $r$ and $\widehat r$,
$|r|-r\sgn(\widehat r)\le2|r-\widehat r|$.
Applying this inequality and then Cauchy--Schwarz across rows,
the expected loss in the unnormalized bilinear sum is at most
\[
 2\sum_i\left(\frac{M}{m_2}
                       \sum_k\mathcal E_{ik}^2\right)^{1/2}
 \le 2M\left(\frac{M}{m_2}B^{-0.20}\right)^{1/2}
 =O(M B^{-0.01}).
\]
Here $m_2\asymp M B^{-0.18}$ for large $B$.  Hence at least
one subset has this loss bound, with
$\sgn(\widehat r_i)=\widetilde g_i$.  Its signs on $J'$ are
the restrictions of the chosen optimizing column signs.

Deleting every ordered pair with an endpoint in $J'$ changes
any such bilinear sum in absolute value by at most
$2m_2d_B$, using symmetry and the degree event.  After deletion,
maximizing the column signs gives exactly
$\sum_{k\in R'}|\sum_{i\in R'}\mathcal E_{ik}\widetilde g_i|$.
Since $m_2d_B/M=O(B^{-0.11})$, the claim follows.
\end{proof}

\begin{lemma}[Uniform mean after fixing a column set]
\label{lem:sampling-conditional-mean}
Fix a deterministic subset $J'$ of size $m_2$ and a deterministic
assignment $\sigma\in\{-1,1\}^{J'}$.  Condition on its site types
$\mathbf S_{J'}=\mathbf t$.  For $i\in R'$, define the fixed
single-site function
\[
 g_i(u)=\sgn\left(\sum_{j\in J'}
                 \mathcal E_{ij}(u,t_j;s)\sigma_j\right).
\]
Then, writing
\[
 Z_{J',\sigma,\mathbf t}
   =\frac1M\sum_{k\in R'}
       \left|\sum_{i\in R'}\mathcal E_{ik}(S_i,S_k;s)g_i(S_i)\right|,
\]
we have
$\E[Z_{J',\sigma,\mathbf t}\mid\mathbf S_{J'}=\mathbf t]
\le E_0+o(1)$, uniformly in $J'$, $\sigma$, $\mathbf t$, and $s$.
\end{lemma}

\begin{proof}
Only the types in the deterministic set $J'$ have been conditioned
on.  In particular, no event describing which signs an optimizer
would choose has been imposed.  The types in $R'$ remain independent
with their usual laws, and $|g_i|\le1$.

For $k\in R'$ and a possible value $v$ of $S_k$, put
\[
 a_k(v)=\sum_{\substack{i\in R'\\ i\ne k}}
             \E_{S_i}\big[\mathcal E_{ik}(S_i,v;s)g_i(S_i)\big].
\]
Choose $h_k(v)=\sgn(a_k(v))$, and set both families of tests
equal to zero on removed indices.  Equation~\eqref{eq:44},
which is uniform over all bounded single-site tests, gives
\[
 \frac1M\sum_{k\in R'}\E_{S_k}|a_k(S_k)|
 =\frac1M\sum_{i,k\in R'}
       \E[\mathcal E_{ik}g_i(S_i)h_k(S_k)]
 \le E_0+o(1).
\]
This application is valid for every fixed value of $\mathbf t$,
however the resulting functions $g_i$ depend on that value.

Conditional on $S_k$, the centered terms
\[
 \mathcal E_{ik}(S_i,S_k;s)g_i(S_i)
  -\E_{S_i}[\mathcal E_{ik}(S_i,S_k;s)g_i(S_i)],
       \qquad i\in R'\setminus\{k\},
\]
are independent.  The diagonal contributes zero.  Conditional
variance followed by Cauchy--Schwarz therefore gives
\[
 \E\left|\sum_{i\in R'}\mathcal E_{ik}g_i(S_i)-a_k(S_k)\right|
 \le\left(\sum_{i\in R'}\E\mathcal E_{ik}^2\right)^{1/2}
 \le q_B^{1/2}.
\]
Symmetry supplies the squared-column bound from the squared-row
bound.  Averaging this estimate over $k$ costs at most
$q_B^{1/2}=O(B^{-0.105})$.  This proves the asserted mean bound,
including its uniformity in the conditioned choice.
\end{proof}

\begin{proof}[Proof of Proposition~\ref{prop:sampling-comparison}]
The skeleton lemma reduces the adaptive cut norm to a finite family of
optimizations, each with conditional mean at most $E_0+o(1)$. The remaining
step is a concentration estimate strong enough to hold for the entire
family. The common degree event will be excluded only once.
We prove concentration for each of the deterministic subset/sign
assignments in Lemma~\ref{lem:sampling-conditional-mean}, and then
take a union over those assignments.  Fix such an assignment and
condition on $\mathbf S_{J'}=\mathbf t$.  Abbreviate its random
quantity by $Z$.  Let $\mathcal G_{R'}$ be the set of configurations
of the remaining site types satisfying
\[
 \max_{i\in R'}\sum_{k\in R'}
                \mathcal H_{ik}(S_i,S_k;s)\le d_B.
\]
The observation after Lemma~\ref{lem:sampling-good} shows that
\[
 \Prob(\mathcal G_{R'}^c\mid\mathbf S_{J'}=\mathbf t)
       =O(B^{-0.03}),
\]
uniformly in the conditioning and assignment.

We next verify a Lipschitz bound between any two configurations
$\mathbf u,\mathbf v\in\mathcal G_{R'}$, rather than only
between good configurations differing in one coordinate.  Let
$A=\{i\in R':u_i\ne v_i\}$.  The row functions $g_i$ have
already been fixed by $\mathbf t$ and $\sigma$.  Thus a term
$\mathcal E_{ik}(u_i,u_k)g_i(u_i)$ can change only if $i$ or $k$
belongs to $A$.  The reverse triangle inequality for each column
sum gives
\[
 |Z(\mathbf u)-Z(\mathbf v)|
 \le\frac1M
   \sum_{\substack{i,k\in R'\\ i\in A\ \text{or}\ k\in A}}
       \big(\mathcal H_{ik}(u_i,u_k;s)
             +\mathcal H_{ik}(v_i,v_k;s)\big)
 \le\frac{4d_B}{M}|A|.
\]
The last inequality uses both endpoint degree bounds and symmetry.
It also includes the change of $g_i$ when $i\in A$.  No connecting
path of good configurations is required.

Choose a fixed $L'>E_0+2+\epsilon_2$.  On $\mathcal G_{R'}$
the function $f=\min(Z,L')$ is nonnegative and Lipschitz for the
site Hamming distance $d_H$, with constant $c_B=4d_B/M$.
For sufficiently large $B$ the good set is nonempty.  We use the
McShane extension \citep{McShane1934} in its infimum form
\citep[Section~3, Equation~(3.2)]{Caputti1984}, followed by clamping:
\[
 \widetilde Z(\mathbf u)
 =\min\left\{L',\ \inf_{\mathbf v\in\mathcal G_{R'}}
             \big(f(\mathbf v)+c_Bd_H(\mathbf u,\mathbf v)\big)
       \right\}.
\]
The site spaces are finite for every $B$, so the formula is
measurable.  The infimum defines a $c_B$-Lipschitz function by
the triangle inequality, and its restriction to the good set is
$f$: the Lipschitz inequality for $f$ gives the lower bound,
and taking $\mathbf v=\mathbf u$ gives the upper bound.
Clamping preserves the Lipschitz bound.  Hence
$0\le\widetilde Z\le L'$ everywhere and
$\widetilde Z=\min(Z,L')$ on $\mathcal G_{R'}$.
By Lemma~\ref{lem:sampling-conditional-mean},
\[
 \E[\widetilde Z\mid\mathbf S_{J'}=\mathbf t]
 \le\E[Z\mid\mathbf S_{J'}=\mathbf t]
      +L'\Prob(\mathcal G_{R'}^c\mid\mathbf S_{J'}=\mathbf t)
 \le E_0+o(1),
\]
again uniformly in every fixed choice.

Changing any one of the independent remaining site types changes
$\widetilde Z$ by at most $c_B$.  McDiarmid's bounded-differences
inequality \citep{McDiarmid1989} therefore gives, for every $a>0$,
\[
 \Prob\left(\widetilde Z-\E[\widetilde Z\mid\mathbf S_{J'}]
                     >a\,\middle|\,\mathbf S_{J'}=\mathbf t\right)
 \le\exp\left(-\frac{2a^2}{|R'|c_B^2}\right)
 \le\exp\left(-\frac{a^2M}{8d_B^2}\right).
\]
Taking $a=\epsilon_2/4$ and using the uniform mean bound shows
that, for large $B$,
\[
 \Prob\left(\widetilde Z>E_0+\epsilon_2/2
                     \,\middle|\,\mathbf S_{J'}=\mathbf t\right)
 \le\exp(-c_{\epsilon_2}M/B^{0.14}).
\]
On $\mathcal G_{R'}$, the event $Z>E_0+\epsilon_2/2$ is
equivalent to this event for $\widetilde Z$, since the threshold
is strictly below $L'$.

There are at most
\[
 N_B=\binom{M}{m_2}2^{m_2}
       \le\left(\frac{2eM}{m_2}\right)^{m_2}
\]
deterministic subset/sign assignments.  For each assignment,
integrate the last conditional tail bound over its subset types.
We may then union-bound the extension tails, because
\[
 \log N_B=O(m_2\log M)=O(B^{0.14}\log B)
       =o(M/B^{0.14}),
 \qquad M/B^{0.14}\asymp B^{0.18}.
\]
The original full-sample event $\Omega_{\rm deg}$ implies
$\mathcal G_{R'}$ for every subset $J'$ simultaneously, by
nonnegativity of the envelope.  Consequently
\[
 \Prob\left(\Omega_{\rm deg}\ \text{and some assignment has}
          \ Z_{J',\sigma,\mathbf S_{J'}}>E_0+\epsilon_2/2\right)
 \le N_B\exp(-c_{\epsilon_2}M/B^{0.14})=o(1).
\]
Only the exponential extension tails have been union-bounded;
the probability of $\Omega_{\rm deg}^c$ is counted once.

With probability $1-o(1)$, the squared-mass event also holds,
and Lemma~\ref{lem:sampling-skeleton} now bounds the fully
adaptive cut norm by
$E_0+\epsilon_2/2+O(B^{-0.01})+O(B^{-0.11})$.
On the exceptional event use $\|\mathcal E\|_\square\le A_B$.
Lemma~\ref{lem:sampling-good} makes its expected contribution
$o(1)$.  This proves \eqref{eq:45}.
All probabilities and errors used here were uniform for each
fixed $s$ in the compact range, which proves the stated
uniformity without constructing a simultaneous event over $s$.
\end{proof}

The deterministic row estimate \eqref{eq:kernel-main-row-bounds} gives
\[
 \sum_{i,k}|\mathcal M_{ik}|
 \ll_{\eta,m_1}M=O(B^{0.32}),
\]
uniformly in the site types and in $s$. Thus $\mathcal M$ satisfies the
$O(B^{10})$ total-mass hypothesis of the independent-root comparison.

Finally, combine \eqref{eq:45} with the independent-site comparison
\eqref{eq:21}.  For the actual sets
$S_i(n)=\{p\in\mathcal P:p\mid n+i\}$ and fixed $B$, the upper
limit as $x\to\infty$ of the absolute error in replacing the kernel
in \eqref{eq:19} by
$\mathcal M_{ik}(S_i(n),S_k(n);n/(Tx))$ is at most
\[
 C_{\rm lab}(E_0+\epsilon_2)+o(1).
\]
Here $x$ tends along the fixed bad subsequence, and the $o(1)$ tends
to zero as $B\to\infty$ after this inner upper limit.  The constant
$C_{\rm lab}$ depends
only on the fixed label bound and the compact $s$-range.
Indeed bounded complex label energies are controlled by a fixed
multiple of the real cut norm.  For each fixed $B$, the error
norm is a finite maximum of finite-residue functions continuous
in $s$ (piecewise continuity would also suffice), so ordinary
scale counting gives its Haar expectation integrated over $s$,
as explained after \eqref{eq:19}.  The uniform expectation bound
above can then be integrated over that fixed compact interval.
This uses no independence between the actual labels and site
types, and all matrices retain zero entries on nonallowed pairs.
